\documentclass[11pt]{article}
\usepackage[T1]{fontenc}
\usepackage{lmodern}
\usepackage[a4paper,margin=25mm]{geometry}
\usepackage{amsmath,amssymb,amsthm}
\usepackage[hypertexnames=false,hidelinks]{hyperref}
\newcommand{\V}{\mathsf V}
\newcommand{\E}{\mathsf E}
\newcommand{\loc}[1]{\par\smallskip\noindent{\small\textit{Source: #1}}\par\smallskip}
\setlength{\emergencystretch}{2em}
\begin{document}
\begin{center}
{\large Probability on Trees and Networks\\[2pt]Exercise 7.10 and its Author Comment}\\[4pt]
Russell Lyons\quad Yuval Peres\\[2pt]
Version of 19 August 2026---corrected and updated paperback edition
\end{center}

% BEGIN CONTEXT 7.10
\noindent\textbf{Source and scope (editorial).}
This private reader reproduces selected English author text from the fixed 716-page paperback PDF (21,010,205 bytes). Exercise 7.10 is on physical p.\ 289 (printed p.\ 272), in the collected in-text exercises of Chapter 7. Its complete Comment is on physical pp.\ 638--639 (printed pp.\ 621--622). The definition of $\rho(G)$ that the exercise quantifies over is printed in the body on physical p.\ 278 (printed p.\ 261), and the two ingredients named by Solution 1, inequality (6.13) and the proof of Proposition 6.6, are printed on physical pp.\ 201--202 (printed pp.\ 184--185). Fekete's lemma, named by Solution 2, is Exercise 3.9, printed on physical p.\ 102 (printed p.\ 85), with its Comment on physical p.\ 622 (printed p.\ 605); the lemma and its Comment are reprinted below as uncounted context because the exercise's own Comment states that the argument uses them.

The separate errata sheet dated 19 August 2026 contains no entry for Chapter 7; accordingly no revision is applied anywhere in this reader, and no printed text here is altered. No source mathematics is repaired or independently certified.

\medskip\noindent\textbf{Use of the exercise in the body of the work.}
\loc{Proposition 7.34 and its proof, physical pp.\ 278--279 (printed pp.\ 261--262).}
The exercise is used as a tool in the proof of Proposition 7.34 (Spectral Radius of Products), twice. On physical p.\ 279 the proof says, for the upper bound, ``For an upper bound, we may use the inequality of Exercise 7.10,'' and, for the lower bound, ``For a lower bound, we note that according to Exercise 7.10, it suffices to consider only even $n$, which we now do.'' Both sentences refer to the two assertions proved in the Comment below.

\subsection*{Definition of the spectral radius}
\loc{Physical p.\ 278 (printed p.\ 261).}
However, although it is crude, quite a useful estimate of $\gamma(G)$ arises from the spectral radius of $G$, which is much easier to calculate. For any graph $G$, recall that the spectral radius of (simple random walk on) $G$ is defined to be
\[
\rho(G):=\limsup_{n\to\infty}\mathbf P_o[X_n=o]^{1/n},
\]
where $\langle X_n\rangle$, $\mathbf P_o$ denotes simple random walk on $G$ starting at $o$.

\subsection*{The two named ingredients: (6.13) and Proposition 6.6}
\loc{Physical pp.\ 201--202 (printed pp.\ 184--185).}
Suppose that $\langle X_n\rangle$ is a Markov chain on a countable state space $\V$ with a stationary measure $\pi$. We define the transition operator
\[
(Pf)(x):=\mathbf E_x f(X_1)=\sum_{y\in\V}p(x,y)f(y).
\]
Then $P$ maps $\ell^2(\V,\pi)$ to itself with norm
\[
\|P\|_\pi:=\|P\|_{\ell^2(\V,\pi)}:=\sup\left\{\frac{\|Pf\|_\pi}{\|f\|_\pi};\ f\ne 0\right\}
\]
at most 1. As the reader should recall from the theory of Markov chains, we have that $(P^nf)(x)=\sum_{y\in\V}p_n(x,y)f(y)$ when $p_n(x,y):=\mathbf P_x[X_n=y]$.

Let $G$ be a (connected) graph with conductances $c(e)>0$ on the edges and $\pi(x)$ be the sum of the conductances incident to a vertex $x$. The operator $P$ that we defined earlier is self-adjoint: For functions $f,g\in D_{00}$, we have
\[
(Pf,g)_\pi=\sum_{x\in\V}\pi(x)(Pf)(x)g(x)=\sum_{x\in\V}\pi(x)\Bigl[\sum_{y\in\V}p(x,y)f(y)\Bigr]g(x)=\sum_{x\in\V}\sum_{y\in\V}c(x,y)f(y)g(x).
\]
Since this is symmetric in $f$ and $g$, it follows that $(Pf,g)_\pi=(f,Pg)_\pi$. Since the functions with finite support are dense in $\ell^2(\V,\pi)$, we get this identity for all $f,g\in\ell^2(\V,\pi)$.

\medskip\noindent\textbf{Proposition 6.6.} \textit{For any two vertices $x,y\in\V$, we have}
\begin{equation}
\|P\|_\pi=\limsup_{n\to\infty}\Bigl(\sup_z p_n(x,z)\Big/\sqrt{\pi(z)}\Bigr)^{1/n}=\limsup_{n\to\infty}p_n(x,y)^{1/n}.\tag{6.12}
\end{equation}
\textit{Moreover,}
\begin{equation}
\forall n\qquad p_n(x,y)\le\sqrt{\pi(y)/\pi(x)}\,\|P\|_\pi^n.\tag{6.13}
\end{equation}
\begin{proof}
Irreducibility of the Markov chain implies that the right-most quantity in (6.12) does not depend on the choice of $x,y$. Thus, define
\begin{equation}
\rho(G):=\limsup_{n\to\infty}p_n(x,y)^{1/n}.\tag{6.14}
\end{equation}
It is clear that the middle term in (6.12) is at least $\rho(G)$. To see that this middle term is at most $\|P\|_\pi$, use $p_n(x,z)=\langle\mathbf 1_{\{x\}},P^n\mathbf 1_{\{z\}}\rangle_\pi\big/\pi(x)$ and the Cauchy--Schwarz inequality to deduce that $p_n(x,z)\le\sqrt{\pi(z)/\pi(x)}\,\|P^n\|_\pi\le\sqrt{\pi(z)/\pi(x)}\,\|P\|_\pi^n$. This also implies (6.13).

Finally, to show the converse inequality that $\|P\|_\pi\le\rho(G)$, suppose that $0\le f\in D_{00}\setminus\{0\}$. Then self-adjointness of $P$ and the Cauchy--Schwarz inequality yield
\[
\|P^{n+1}f\|_\pi^4=(P^nf,P^{n+2}f)_\pi^2\le\|P^nf\|_\pi^2\|P^{n+2}f\|_\pi^2.
\]
This means that $\|P^{n+1}f\|_\pi/\|P^nf\|_\pi$ is increasing, whence it has a limit, $L$. Since the product of these quotients for $n=0,\dots,N-1$ is $\|P^Nf\|_\pi/\|f\|_\pi$, we have that
\begin{equation}
L=\lim_{N\to\infty}\bigl(\|P^Nf\|_\pi/\|f\|_\pi\bigr)^{1/N}=\lim_{N\to\infty}\|P^Nf\|_\pi^{1/N}.\tag{6.15}
\end{equation}
Now $p_{2n}(x,x)\ge p_n(x,x)^2$, since one way to return after $2n$ steps is to return after $n$ steps and then again after $n$ steps. Therefore,
\[
\rho(G)=\limsup_{n\to\infty}p_{2n}(x,x)^{1/2n}.
\]
Thus, for $0\le f\in D_{00}\setminus\{0\}$, we have, by self-adjointness of $P$,
\[
\limsup_{n\to\infty}\|P^nf\|_\pi^{1/n}=\limsup_{n\to\infty}(P^{2n}f,f)_\pi^{1/2n}=\limsup_{n\to\infty}\Bigl[\sum_{x,y}\pi(x)f(x)f(y)p_{2n}(x,y)\Bigr]^{1/2n}=\rho(G).
\]
Combining this with (6.15), we deduce that $L=\rho(G)$. In particular, $\|Pf\|_\pi/\|f\|_\pi\le\rho(G)$. When $f\in D_{00}$ is not assumed to be nonnegative, we may use the inequality $\|Pf\|_\pi\le\|P|f|\|_\pi$ to deduce that $\|Pf\|_\pi\le\rho(G)\|f\|_\pi$ in general, in other words, $\|P\|_\pi\le\rho(G)$.
\end{proof}

\noindent The book's footnote to this conclusion reads: ``For a general operator on a Banach space, its spectral radius is defined to be $\max|z|$ for $z$ in the spectrum of the operator. For reversible Markov operators acting on $\ell^2$, this agrees with the definition we gave; see, for example, Taylor and Lay (1980), Theorem VI.3.3. We will not need this general representation explicitly. For a nonreversible Markov chain, the norm of the transition operator is not generally equal to its spectral radius.''

\subsection*{Fekete's lemma: Exercise 3.9 and its Comment}
\loc{Physical p.\ 102 (printed p.\ 85).}
One very nice feature of subperiodic trees is that this equality also holds for them! This fact, as well as its proof, is analogous to a classical fact about sequences of reals, known as Fekete's lemma (Exercise 3.9). A sequence $\langle a_n\rangle$ of real numbers is called \emph{subadditive} if
\[
\forall m,n\ge1\qquad a_{m+n}\le a_m+a_n.
\]
A simple example is $a_n:=\lceil\beta n\rceil$ for some real $\beta>0$.

\medskip\noindent\textbf{Exercise 3.9.}
(a) (Fekete's Lemma) Show that for every subadditive sequence $\langle a_n\rangle$, the sequence $\langle a_n/n\rangle$ converges to its infimum:
\[
\lim_{n\to\infty}\frac{a_n}{n}=\inf_n\frac{a_n}{n}.
\]
(b) Show that Fekete's lemma holds even if a finite number of the $a_n$ are infinite. (c) Show that for every $0$-subperiodic tree $T$, the limit $\lim_{n\to\infty}|T_n|^{1/n}$ exists.

\loc{Comment on Exercise 3.9, physical p.\ 622 (printed p.\ 605).}
3.9. (a) Let $\beta>\inf a_n/n$ and let $m$ be such that $a_m/m<\beta$. Write $a_0:=0$. For any $n$, write $n=qm+r$ with $0\le r\le m-1$; then we have
\[
a_n=a_{qm+r}\le a_m+\cdots+a_m+a_r=qa_m+a_r,
\]
whence
\[
\frac{a_n}{n}\le\frac{qa_m+a_r}{qm+r}=\frac{a_m}{m}\cdot\frac{qm}{qm+r}+\frac{a_r}{n}<\beta\cdot\frac{qm}{qm+r}+\frac{a_r}{n},
\]
and so
\[
\limsup_{n\to\infty}\frac{a_n}{n}\le\beta.
\]
(b) Modify the proof so that if an infinite $a_r$ appears, then it is replaced by $a_{m+r}$ instead. (c) Observe that $\log|T_n|$ is subadditive.
% END CONTEXT 7.10

\clearpage
% BEGIN PRIMARY STATEMENT 7.10
\section*{Exercise 7.10}
\loc{Physical p.\ 289 (printed p.\ 272).}
\begin{sloppypar}
\noindent\textbf{Exercise 7.10.} Show that for any graph $G$, the spectral radius satisfies $\rho(G)=\lim_{n\to\infty}\mathbf P_o[X_{2n}=o]^{1/2n}$ and $\mathbf P_o[X_n=o]\le\rho(G)^n$ for all $n$.
\end{sloppypar}
% END PRIMARY STATEMENT 7.10

% BEGIN PRIMARY PROOF 7.10
\section*{Complete author Comment on Exercise 7.10}
\loc{Physical pp.\ 638--639 (printed pp.\ 621--622).}
\textbf{7.10.} Solution 1. This follows from (6.13) and the proof of Proposition 6.6.

Solution 2. A direct proof goes as follows. Let us first suppose that $\mathbf P_o[X_k=o]$ is strictly positive for all large $k$. Note that for any $k$ and $n$, we have
\[
\mathbf P_o[X_n=o]\ge\mathbf P_o[X_k=o]\,\mathbf P_o[X_{n-k}=o].
\]
Thus, by Fekete's lemma (Exercise 3.9), we have $\rho(G)=\lim_{n\to\infty}\mathbf P_o[X_n=o]^{1/n}$ as well as $\mathbf P_o[X_n=o]\le\rho(G)^n$ for all $n$.

On the other hand, if $\mathbf P_o[X_k=o]=0$ for odd $k$, it is still true that $\mathbf P_o[X_k=o]$ is strictly positive for all even $k$, whence $\lim_{n\to\infty}\mathbf P_o[X_{2n}=o]^{1/2n}$ exists and equals $\rho(G)$.
% END PRIMARY PROOF 7.10

\medskip\noindent\textbf{Cached bibliographic attribution (not an additional source proof).}
The Comment on Exercise 7.10 carries no external bibliographic attribution. Its two named inputs are internal to the book: printed inequality (6.13) with Proposition 6.6 (physical pp.\ 201--202), and Fekete's lemma as Exercise 3.9 (physical p.\ 102, Comment physical p.\ 622), both reprinted above as uncounted context.
\end{document}
